\documentclass[12pt]{article}
\usepackage[utf8]{inputenc}
\usepackage[italian]{babel}
\usepackage{amsmath, amssymb, amsfonts}
\usepackage{tikz}
\usepackage{geometry}
\geometry{a4paper, margin=2cm}

\begin{document}

\section*{TEOREMA DI POINSOT ED ELLISSOIDE DI INERZIA}

\begin{minipage}{0.3\textwidth}
\begin{tikzpicture}
\draw[->] (0,0) -- (2,0) node[right] {$y$};
\draw[->] (0,0) -- (0,2) node[above] {$z$};
\draw[->] (0,0) -- (-1.5,-1) node[below left] {$x$};
\draw[->, thick, olive] (0,0) -- (1.5,1.5) node[right] {$\hat{u}$};
\draw[->] (0,0) -- (0.8,2) node[right] {$dm$};
\draw[dashed] (0.8,2) -- (0.56, 0.56) node[midway, right] {$R$};
\node at (0.4, 0.9) {$\phi$};
\node at (0.3, 1.3) {$\vec{r}$};
\end{tikzpicture}
\end{minipage}
\begin{minipage}{0.65\textwidth}
$\hat{u} = $ asse di rotazione, definito da $u_x, u_y, u_z$ \\
$u_x^2 + u_y^2 + u_z^2 = 1$ \\[0.2cm]
$R = ||\vec{r}|| \sin\phi = ||\vec{r} \times \hat{u}|| = $ distanza fra $dm$ e $u$ \\[0.2cm]
$I_u = \int_{CR} dm ||\vec{r} \times \hat{u}||^2 = \int_{CR} dm \left[ (y u_z - z u_y)^2 + (z u_x - x u_z)^2 + (x u_y - y u_x)^2 \right] = \sum_{i,j} \int_{CR} dm (i u_j - j u_i)^2$
\end{minipage}

\vspace{0.5cm}
\noindent $\Rightarrow I_u = u_x^2 I_{xx} + u_y^2 I_{yy} + u_z^2 I_{zz} + 2 u_x u_y I_{xy} + 2 u_x u_z I_{xz} + 2 u_y u_z I_{yz}$

\vspace{0.5cm}
\noindent Se dividiamo tutto per $I_u$ e consideriamo \\
$A = \frac{u_x}{\sqrt{I_u}}, \quad B = \frac{u_y}{\sqrt{I_u}}, \quad C = \frac{u_z}{\sqrt{I_u}} \quad$ allora si ottiene:

\vspace{0.3cm}
\fbox{$I_{xx} A^2 + I_{yy} B^2 + I_{zz} C^2 + 2(AB I_{xy} + BC I_{yz} + AC I_{xz}) = 1$} \\
= ELLISSOIDE DI INERZIA

\vspace{0.3cm}
\begin{minipage}{0.65\textwidth}
$P = (A,B,C) \quad OP = \sqrt{A^2 + B^2 + C^2} = \frac{1}{\sqrt{I_u}} \sqrt{u_x^2 + u_y^2 + u_z^2}^{=1}$ \\[0.2cm]
\fbox{$\Rightarrow I_u = \frac{1}{OP^2}$} \quad rispetto a una terna di assi qualsiasi l'ellissoide ha formula quadratica arbitraria $(A,B,C)$
\end{minipage}
\begin{minipage}{0.3\textwidth}
\begin{tikzpicture}
\draw[->] (0,0) -- (2,0) node[right] {$B$};
\draw[->] (0,0) -- (0,1.5) node[above] {$C$};
\draw[->] (0,0) -- (-1,-1) node[below left] {$A$};
\draw[olive, thick] (0,0) ellipse (1.5 and 0.8);
\draw[dashed, ->] (0,0) -- (1.2, 0.5) node[right] {$P$};
\draw[dashed] (0,0) -- (1.5, -0.4) node[right] {$u$};
\end{tikzpicture}
\end{minipage}

\vspace{0.5cm}
\begin{itemize}
    \item Se cambiamo asse e consideriamo $u'$ (e quindi $I_{u'}$) sicuramente l'equazione dell'ellissoide sarà diversa ma non cambia $I_{u'} = 1 / \overline{OP}^2$ cioè il significato.
    \item[$\Rightarrow$] L'ellissoide è una MAPPATURA della GEOMETRIA del corpo in un certo SISTEMA DI ASSI: se cambio gli assi $I_{ij}, A, B, C$ cambiano ma l'ellissoide, essendo una proprietà del corpo, rimane lo stesso.
\end{itemize}

\vspace{0.3cm}
\noindent Nel caso in cui $x^P, y^P, z^P$ siano principali (quindi $x, y, z$ originari) allora $I_{ii} = I_{ii}^P$ e $I_{ij} = 0 \ \forall i,j$ quindi
\fbox{$I_{xx}^P A^2 + I_{yy}^P B^2 + I_{zz}^P C^2 = 1$} è l'ellissoide rispetto ad assi principali e in questo caso i semiassi sono $(a,b,c)$: \\
$a = 1/\sqrt{I_{xx}^P}, \quad b = 1/\sqrt{I_{yy}^P}, \quad c = 1/\sqrt{I_{zz}^P}$

\subsection*{Esempio: SFERA E CILINDRO}

\begin{minipage}{0.5\textwidth}
\begin{tikzpicture}
\draw[->] (0,0) -- (2,0) node[right] {$B$};
\draw[->] (0,0) -- (0,2) node[above] {$C$};
\draw[->] (0,0) -- (-1,-1) node[below left] {$A$};
\draw[olive, thick] (0,0) circle (1.2);
\draw[dashed] (0,0) -- (1.2,0) node[midway, below] {$b$};
\draw[dashed] (0,0) -- (0,1.2) node[midway, left] {$c$};
\draw[dashed] (0,0) -- (-0.84,-0.84) node[midway, above left] {$a$};
\end{tikzpicture}

\textbf{SFERA}
\begin{itemize}
    \item $\infty$ assi $x^P, y^P, z^P$
    \item $I_{xx}^P = I_{yy}^P = I_{zz}^P = I^P = \frac{2}{5}MR^2$
\end{itemize}
= SFERA di raggio \\
\fbox{$A^2 + B^2 + C^2 = \frac{5}{2} \frac{1}{MR^2}$} $\rightarrow a=b=c$
\end{minipage}
\begin{minipage}{0.45\textwidth}
\begin{tikzpicture}
\draw[->] (0,0) -- (2.5,0) node[right] {$B$};
\draw[->] (0,0) -- (0,1.5) node[above] {$C$};
\draw[->] (0,0) -- (-1,-1) node[below left] {$A$};
\draw[olive, thick] (0,0) ellipse (2 and 0.6);
\draw[dashed] (0,0) -- (2,0) node[midway, below] {$b$};
\draw[dashed] (0,0) -- (0,0.6) node[midway, left] {$c$};
\draw[dashed] (0,0) -- (-0.42,-0.42) node[midway, above left] {$a$};
\end{tikzpicture}

\textbf{CILINDRO}
\begin{itemize}
    \item considero $x^P, y^P, z^P$
    \item $I_{xx}^P = I_{zz}^P = \frac{1}{4}MR^2 + \frac{1}{12}ML^2$ \\
          e \\
          $I_{yy}^P = \frac{1}{2}MR^2$
\end{itemize}
= ELLISSOIDE CON ASSI DI RIVOLUZIONE ($a=c, b$)
\end{minipage}

\end{document}
